\documentclass[10pt,a4paper]{amsart}
\usepackage{amsmath,amssymb,amsthm,graphicx}
\usepackage{verbatim}
\newenvironment{reference}{\comment}{\endcomment}
%\newenvironment{reference}{}{}
\newenvironment{slogan}{\comment}{\endcomment}
\newenvironment{history}{\comment}{\endcomment}
% For commutative diagrams we use Xy-pic
\usepackage[all]{xy}

% We use 2cell for 2-commutative diagrams.
\xyoption{2cell}
% Theorem environments.
%
\theoremstyle{plain}
\newtheorem{theorem}[subsection]{Theorem}
\newtheorem{proposition}[subsection]{Proposition}
\newtheorem{lemma}[subsection]{Lemma}

\theoremstyle{definition}
\newtheorem{definition}[subsection]{Definition}
\newtheorem{example}[subsection]{Example}
\newtheorem{exercise}[subsection]{Exercise}
\newtheorem{situation}[subsection]{Situation}

\theoremstyle{remark}
\newtheorem{remark}[subsection]{Remark}
\newtheorem{remarks}[subsection]{Remarks}

\numberwithin{equation}{subsection}

% Macros
%
\def\lim{\mathop{\mathrm{lim}}\nolimits}
\def\colim{\mathop{\mathrm{colim}}\nolimits}
\def\Spec{\mathop{\mathrm{Spec}}}
\def\Hom{\mathop{\mathrm{Hom}}\nolimits}
\def\Ext{\mathop{\mathrm{Ext}}\nolimits}
\def\SheafHom{\mathop{\mathcal{H}\!\mathit{om}}\nolimits}
\def\SheafExt{\mathop{\mathcal{E}\!\mathit{xt}}\nolimits}
\def\Sch{\mathit{Sch}}
\def\Mor{\mathop{\mathrm{Mor}}\nolimits}
\def\Ob{\mathop{\mathrm{Ob}}\nolimits}
\def\Sh{\mathop{\mathit{Sh}}\nolimits}
\def\NL{\mathop{N\!L}\nolimits}
\def\CH{\mathop{\mathrm{CH}}\nolimits}
\def\proetale{{pro\text{-}\acute{e}tale}}
\def\etale{{\acute{e}tale}}
\def\QCoh{\mathit{QCoh}}
\def\Ker{\mathop{\mathrm{Ker}}}
\def\Im{\mathop{\mathrm{Im}}}
\def\Coker{\mathop{\mathrm{Coker}}}
\def\Coim{\mathop{\mathrm{Coim}}}

% Boxtimes
%
\DeclareMathSymbol{\boxtimes}{\mathbin}{AMSa}{"02}

%
% Macros for moduli stacks/spaces
%
\def\QCohstack{\mathcal{QC}\!\mathit{oh}}
\def\Cohstack{\mathcal{C}\!\mathit{oh}}
\def\Spacesstack{\mathcal{S}\!\mathit{paces}}
\def\Quotfunctor{\mathrm{Quot}}
\def\Hilbfunctor{\mathrm{Hilb}}
\def\Curvesstack{\mathcal{C}\!\mathit{urves}}
\def\Polarizedstack{\mathcal{P}\!\mathit{olarized}}
\def\Complexesstack{\mathcal{C}\!\mathit{omplexes}}
% \Pic is the operator that assigns to X its picard group, usage \Pic(X)
% \Picardstack_{X/B} denotes the Picard stack of X over B
% \Picardfunctor_{X/B} denotes the Picard functor of X over B
\def\Pic{\mathop{\mathrm{Pic}}\nolimits}
\def\Picardstack{\mathcal{P}\!\mathit{ic}}
\def\Picardfunctor{\mathrm{Pic}}
\def\Deformationcategory{\mathcal{D}\!\mathit{ef}}
\usepackage{cleveref}
\crefname{originaldefinition}{original definition}{original definitions}
\Crefname{originaldefinition}{Original definition}{Original definitions}
\makeatletter
\newlabel{definition-chern-classes}{{02U0}{source}{}{sourcepointer-1}{}}
\newlabel{definition-chern-classes@cref}{{[originaldefinition][0][]02U0}{[1][0][]source}{}{}{}}
\makeatother
\crefname{originaldefinition}{original definition}{original definitions}
\Crefname{originaldefinition}{Original definition}{Original definitions}
\makeatletter
\newlabel{definition-divisor-invertible-sheaf}{{02SJ}{source}{}{sourcepointer-2}{}}
\newlabel{definition-divisor-invertible-sheaf@cref}{{[originaldefinition][0][]02SJ}{[1][0][]source}{}{}{}}
\makeatother
\crefname{originalequation}{original equation}{original equations}
\Crefname{originalequation}{Original equation}{Original equations}
\makeatletter
\newlabel{equation-chern-classes}{{02U1}{source}{}{sourcepointer-3}{}}
\newlabel{equation-chern-classes@cref}{{[originalequation][0][]02U1}{[1][0][]source}{}{}{}}
\makeatother
\crefname{originalreference}{original reference}{original references}
\Crefname{originalreference}{Original reference}{Original references}
\makeatletter
\newlabel{situation-setup}{{02QL}{source}{}{sourcepointer-4}{}}
\newlabel{situation-setup@cref}{{[originalreference][0][]02QL}{[1][0][]source}{}{}{}}
\makeatother
\title{On an integral scheme the first Chern class coincides with the Weil divisor associated to the invertible sheaf}
\author{The Stacks Project Authors}
\date{stacks-project:chow.tex:local-snapshot-20261012; GNU FDL 1.2 or later}
\begin{document}
\maketitle
\noindent\textit{Editorial note.} All delivered mathematical text is the author's own, copied from the pinned Stacks Project chapter file chow.tex; no statement or proof step was rewritten, reordered or removed.\par\medskip
\noindent\textit{Editorial note.} Preamble provenance: the pinned chapter file opens with \textbackslash{}input\{preamble\} and carries no preamble text of its own, so the project's single Stacks-wide preamble staged at sources/stacks-chow/preamble.tex is excerpted verbatim as three spans, lines 8-12 (the verbatim package together with the project's own reference, slogan and history display-environment definitions, including the author's commented-out older form of the reference environment), lines 14-18 (Xy-pic setup) and lines 154-217 (the native theorem declarations and mathematical macros). The lines 8-12 span is required because a counted statement or its surrounding text may use those display environments, which the project typesets as comments. Nothing was added, removed or redefined; the wrapper class and the amsmath/amssymb/amsthm/graphicx base packages are supplied by the assembler.\par\medskip
\noindent\textit{Editorial note.} Same-chapter pointers are located at their \textbackslash{}label lines inside the pinned chapter file; cross-chapter pointers are located in the project's own tag list staged in the snapshot (sources/stacks-chow/tags.txt), at the verbatim lines carrying their canonical tags.\par\medskip
\noindent\textit{Allowed standard prerequisites (editorial).}\begin{itemize}
\item The defining relation of the Chern classes, obtained by expanding the cycle class of the zero scheme of a section of the tautological quotient. Same chapter chow.tex, Definition \textbackslash{}ref\{definition-chern-classes\} and the equation labelled equation-chern-classes, used at proof line 7200. Specialized to the case of rank one, where \$\textbackslash{}mathbf\{P\}(\textbackslash{}mathcal\{L\}) = X\$ and \$\textbackslash{}mathcal\{O\}\_\{\textbackslash{}mathbf\{P\}(\textbackslash{}mathcal\{L\})\}(1) = \textbackslash{}mathcal\{L\}\$, it reduces to \$(-1)\textasciicircum{}1 c\_1(\textbackslash{}mathcal\{L\}) \textbackslash{}cap c\_0 + (-1)\textasciicircum{}0 c\_1 = 0\$, which gives the claim.
\item The Weil divisor associated to an invertible sheaf. Same chapter chow.tex, Definition \textbackslash{}ref\{definition-divisor-invertible-sheaf\} (invoked at proof line 7196). It supplies the construction that the proof denotes by \$c\_1(\textbackslash{}mathcal\{L\}) \textbackslash{}cap [X]\$ and compares with the first Chern class.
\end{itemize}\medskip
\ifdefined\hypertarget\hypertarget{sourcepointer-1}{}\fi
\noindent\textit{Source reference.} Original definition 02U0: sources/stacks-chow/pinned-original-source.tex, source lines 7144--7144. Navigation pointer for \textbackslash{}ref\{definition-chern-classes\} used inside the counted span (line 7189); the argument that consumes it is retained in the span.\par\medskip
\ifdefined\hypertarget\hypertarget{sourcepointer-2}{}\fi
\noindent\textit{Source reference.} Original definition 02SJ: sources/stacks-chow/pinned-original-source.tex, source lines 4149--4149. Navigation pointer for \textbackslash{}ref\{definition-divisor-invertible-sheaf\} used inside the counted span (line 7191); the argument that consumes it is retained in the span.\par\medskip
\ifdefined\hypertarget\hypertarget{sourcepointer-3}{}\fi
\noindent\textit{Source reference.} Original equation 02U1: sources/stacks-chow/pinned-original-source.tex, source lines 7156--7156. Navigation pointer for \textbackslash{}ref\{equation-chern-classes\} used inside the counted span (line 7200); the argument that consumes it is retained in the span.\par\medskip
\ifdefined\hypertarget\hypertarget{sourcepointer-4}{}\fi
\noindent\textit{Source reference.} Original reference 02QL: sources/stacks-chow/pinned-original-source.tex, source lines 1372--1372. Navigation pointer for \textbackslash{}ref\{situation-setup\} used inside the counted span (line 7184); the argument that consumes it is retained in the span.\par\medskip
\begin{lemma}
\label{lemma-first-chern-class}
Let $(S, \delta)$ be as in Situation \ref{situation-setup}.
Let $X$ be locally of finite type over $S$.
Assume $X$ is integral and $n = \dim_\delta(X)$.
Let $\mathcal{L}$ be an invertible $\mathcal{O}_X$-module.
The first Chern class of $\mathcal{L}$ on $X$ of
Definition \ref{definition-chern-classes}
is equal to the Weil divisor associated to $\mathcal{L}$
by Definition \ref{definition-divisor-invertible-sheaf}.
\end{lemma}

\begin{proof}
In this proof we use $c_1(\mathcal{L}) \cap [X]$ to denote the
construction of Definition \ref{definition-divisor-invertible-sheaf}.
Since $\mathcal{L}$ has rank $1$ we have
$\mathbf{P}(\mathcal{L}) = X$ and
$\mathcal{O}_{\mathbf{P}(\mathcal{L})}(1) = \mathcal{L}$
by our normalizations. Hence (\ref{equation-chern-classes})
reads
$$
(-1)^1 c_1(\mathcal{L}) \cap c_0 + (-1)^0 c_1 = 0
$$
Since $c_0 = [X]$, we conclude $c_1 = c_1(\mathcal{L}) \cap [X]$
as desired.
\end{proof}

\end{document}
